\section{Computations for the noise certificates}\label{sec:ppt3-app-noise}

This appendix contains the proofs of the statements of Section~\ref{sec:ppt3-noise} that are not given there. The rational certificate data
themselves are ancillary (Appendix~\ref{sec:ppt3-verifiers}).

\begin{proof}[Proof of Lemma~\ref{lem:ppt3-dps}]
For $x=a\otimes b$ write $a\otimes a\otimes b=Vc$. Partial transposition of the product projector gives
\[
\|a\|^2\langle x|Y|x\rangle=c\dg Pc+\langle a a\bar b|Z_1|a a\bar b\rangle+\langle a\bar a b|Z_2|a\bar a b\rangle\ \ge0 .
\]
One certificate per orbit of cuts under the symmetries of $Y$ suffices, and cyclic invariance covers all cuts on
three qubits; thus $Y$ is nonnegative on every pure state that is a product across some cut, hence on the convex hull
$\mathrm{BS}$. The noise threshold is the condition $-(1-p)g+ph<0$. The second bound follows from
$|\tr(Y(\sigma-\rho))|\le(u-l)D(\sigma,\rho)$ for $\sigma\in\mathrm{BS}$. For the last assertion, fix a cut and a unit vector
$\phi$, and set $A_\phi=(|\phi\rangle\langle\phi|)^{T_c}$.
If $\lambda_1\ge\lambda_2\ge\cdots$ are the squared
Schmidt coefficients of $\phi$ at that cut (padded by zeros if necessary), the spectrum
of $A_\phi$ consists of $\lambda_i$,
$\pm\sqrt{\lambda_i\lambda_j}$ for $i<j$, and zeros.
Since $\lambda_2\le1-\lambda_1$, its spectral width is bounded by
\[
\lambda_{\max}(A_\phi)-\lambda_{\min}(A_\phi)
=\lambda_1+\sqrt{\lambda_1\lambda_2}
\le\max_{0\le x\le1}\left[x+\sqrt{x(1-x)}\right]
=\frac{1+\sqrt2}{2}.
\]
For $\Delta=\sigma-\rho$ with $\tr\Delta=0$ and
$D(\rho,\sigma)=r$, the positive and negative parts
of $\Delta$ each have trace $r$. Consequently
$|\tr(A_\phi\Delta)|$ is at most $r$ times the
spectral width of $A_\phi$. Since $\phi$ is arbitrary,
\[
 \langle\phi|\sigma^{T_c}|\phi\rangle
 \ge c-\frac{1+\sqrt2}{2}\,r>0.
\]
The same reasoning applies to every cut.
\end{proof}

\begin{proof}[Proof of Corollary~\ref{cor:ppt3-noise}]
Apply Lemma~\ref{lem:ppt3-dps}. The exact verification checks $Z_1,Z_2,P\succ0$, cyclic symmetry, all $144$
coefficients of the product-vector identity, outward spectral bounds and every number in the table. For local noise a
check at the single rate $q_0$ suffices: $\mathcal D_{q_0}$ on all sites is obtained from any rates $q_j\le q_0$ by further
local depolarisation of rates $(q_0-q_j)/(1-q_j)$, which preserves $\mathrm{BS}$, so a GME output at $q_0$ forces a GME
input. For uniformity in $t$, the numerator of $\rho_t$ is affine in $t$, so $7(t+3)\tr(Y\rho_t)$ is affine in $t$, and so are
$7(t+3)\bigl[(1-p)\tr(Y\rho_t)+ph\bigr]$ and $7(t+3)\bigl[\tr(Y\rho_t)+r(u-l)\bigr]$ for the fixed $Y$, with $h=\tr Y/8$,
$\operatorname{spec}Y\subset[l,u]$ and fixed $p$ and ball radius $r$. An affine function that is negative at both endpoints of the
rational bracket $[3.2567,3.3886]$, which contains \eqref{eq:ppt3-interval}, is negative throughout it, and $7(t+3)>0$. By
Lemma~\ref{lem:ppt3-dps} (the noise threshold $-(1-p)g+ph<0$ and the distance bound $D(\rho,\mathrm{BS})\ge g/(u-l)$), the two
endpoint checks therefore certify the stated interval-wide white-noise and GME-ball bounds.
\end{proof}

\begin{proof}[Proof of Corollary~\ref{cor:ppt3-noise-reused}]
Write $Y=Y_*$ for the cyclic rational witness of the
certificate data. Its exact level-two certificate proves $\tr(Y\sigma)\ge0$ on
all biseparable states. The Pythagorean identity
$312^2+266^2=410^2$ makes the coefficients of
$\rho_*$ rational. The exact certificate evaluation gives
\[
  g:=-\tr(Y\rho_*)>\frac{173}{40000},\qquad
  h:=\frac{\tr Y}{8}=1,\qquad
  -\frac14 I\prec Y\prec\frac{31}{10}I .
\]
Thus the spectral width of $Y$ is less than $67/20$.
These inequalities give the white-noise and trace-distance bounds.

For (i), at $p=2153/500000$,
\[
 (1-p)\frac{173}{40000}-p
 =\frac{7531}{20000000000}>0,
\]
so $\tr[Y((1-p)\rho_*+pI/8)]<0$. The inequality
also holds for smaller $p$. Part (ii) is the
exact local-noise check at $q=81/50000$:
all smaller, unequal local rates can be followed by
additional local depolarisation to reach the checked
uniform output. Local channels preserve biseparability,
and these channels preserve full PPT.

For (iii), let $p_0=32/10000$ and $r_0=331/1000000$.
The witness margin remains strictly negative throughout
the trace-distance ball because
\[
 (1-p_0)\frac{173}{40000}-p_0
           -\frac{67}{20}r_0
  =\frac{231}{100000000}>0.
\]
Since $\rho_*^{T_c}\succeq0$, mixing with white
noise also gives $\rho_\circ^{T_c}\succeq p_0I/8$.
The PPT stability estimate in
Lemma~\ref{lem:ppt3-dps} applies because
\[
 \frac{p_0}{8}-\frac{1+\sqrt2}{2}r_0
 =\frac{469-331\sqrt2}{2000000}>0.
\]
Finally, (iv) follows directly from
\[
 \frac{173}{40000}
 -\frac{67}{20}\frac{129}{100000}
 =\frac7{2000000}>0.
\]
The rational inequalities, including the enclosure
$\sqrt2<99/70$, are checked exactly from the certified data
(Appendix~\ref{sec:ppt3-verifiers}).
\end{proof}

We now give the state $\tilde\rho$ explicitly. With $v,w,e,o$ as in \eqref{eq:ppt3-vectors}, write $\ket g=\ket{000}$ and
$\ket h=\ket{011}+\ket{101}+\ket{110}$. Then
\begin{equation}\label{eq:ppt3-tilde-rho}
\begin{aligned}
10^7\,\tilde\rho={}&955592\,\ket g\bra g-930700\,(\ket g\bra h+\ket h\bra g)+906857\,\ket h\bra h\\
&+519407\,\ket e\bra e+208476\,\ket{\bar e}\bra{\bar e}\\
&+390647\,\ket v\bra v-520181\,(\ket v\bra w+\ket w\bra v)+693325\,\ket w\bra w\\
&+556418\,\ket o\bra o+104\,\ket{\bar o}\bra{\bar o}.
\end{aligned}
\end{equation}
These ten coefficients are the state's coordinates in the invariant basis of the cyclic, parity and $K\circ S_{BC}$
symmetries; the witness data of Corollary~\ref{cor:ppt3-robust} are ancillary.

\begin{proof}[Proof of Corollary~\ref{cor:ppt3-robust}]
Lemma~\ref{lem:ppt3-dps} applies. The exact verification checks trace one, full rank, cyclic and parity invariance, all
three partial-transpose eigenvalue lower bounds, and the witness and tolerance inequalities. The PPT-only ball radii here
and in the table of Corollary~\ref{cor:ppt3-noise} follow from those eigenvalue bounds and the stability estimate of
Lemma~\ref{lem:ppt3-dps}; their rational comparisons are checked against the same constants.
\end{proof}

\begin{proof}[Proof of Proposition~\ref{prop:ppt3-noise-sharp}]
The lower bounds are Corollaries~\ref{cor:ppt3-noise-reused}, \ref{cor:ppt3-noise} and~\ref{cor:ppt3-robust}. For the
upper bounds, let $\mathcal C(X)=\frac13\sum_{k=0}^2\Pi^kX\Pi^{-k}$ average over the cyclic qubit permutation $\Pi$. The
certificate data give, for each of the three states, three orbit representatives: a Gaussian-rational qubit vector $a_r$ and a
positive definite Gaussian-rational $4\times4$ matrix $\Lambda_r$. The states are invariant under global parity $Z^{\otimes3}$ and
the antiunitary $K\circ S_{BC}$, which preserve the cut $A|BC$; applying them to $|a_r\rangle\langle a_r|\otimes\Lambda_r$ gives at most four terms
per representative, and $\theta$ is the sum of all of them (at most twelve terms). Thus $\theta=\sum_i|a_i\rangle\langle a_i|\otimes\Lambda_i$ with $\Lambda_i\succ0$
is separable across $A|BC$.
At $p$ equal to the stated upper bound, the residual $\Delta=(1-p)\rho+pI/8-\mathcal C(\theta)$ is cyclically invariant and, with $c=\tr\Delta/8>0$,
satisfies $\|\Delta-cI\|_2\le c$ exactly. By the bipartite Gurvits--Barnum bound~\cite{GB}, this makes $\Delta$
separable across $A|BC$. Hence
$(1-p)\rho+pI/8=\mathcal C(\theta+\Delta)$ is an equal mixture of states separable across the three cuts, so it is biseparable at
$p$ equal to the stated upper bound. For every larger $p'$, $(1-p')\rho+p'I/8$ is a convex mixture of the state at $p$ with the
separable state $I/8$, hence biseparable too.
\end{proof}
